\documentclass[twocolumn,10pt]{article}

% --- Self-contained two-column preprint (no external arxiv.sty needed) --------
\usepackage[utf8]{inputenc}
\usepackage[T1]{fontenc}
\usepackage{lmodern}
\usepackage[margin=0.85in]{geometry}
\usepackage{graphicx}
\usepackage{booktabs}
\usepackage{amsmath}
\usepackage[numbers,sort&compress]{natbib}
\usepackage{xcolor}
\usepackage[colorlinks=true,allcolors=blue]{hyperref}

\graphicspath{{experiments/runs/}{figures/}}

% Numbers/claims to fill in after running experiments are wrapped so they are
% easy to find and visually obvious in a draft build.
\newcommand{\ph}[1]{\textcolor{red}{\textbf{[#1]}}}

% Include an atlas panel if it has been generated, else a labeled placeholder box
% so the paper compiles before experiments are run.
\newcommand{\atlas}[1]{%
  \IfFileExists{experiments/runs/#1.png}%
    {\includegraphics[width=0.32\textwidth]{experiments/runs/#1.png}}%
    {\fbox{\parbox[c][3.4cm][c]{0.30\textwidth}{\centering\footnotesize
       \textcolor{red}{\texttt{\detokenize{#1}}}\\[2pt](panel not generated)}}}%
}

\title{\bf Mapping the Narrow Corridor with Large Language Models}

\author{Ebrahim M. Songhori \\ \texttt{e.songhori@gmail.com}}
\date{July 1, 2026}

\begin{document}
\maketitle

\begin{abstract}
Acemoglu and Robinson's \emph{The Narrow Corridor} traces national histories as
paths through a two-dimensional (state power, society power) space but plots
none: quantifying the two axes at each moment is the labor-intensive,
judgment-laden task that has kept the framework qualitative. We ask whether large
language models (LLMs) can operationalize it. We introduce a reproducible,
provider-agnostic pipeline that scores a country period by period using
chain-of-thought (events first, then score) and in-context anchoring on prior
periods, against an explicit rubric with schema-validated output. As a proof of
concept we produce a six-country trajectory atlas (Iran, France, the United
Kingdom, the United States, China, Chile) and compare four models. We assess the
scores for consistency with an expert index (V-Dem) and for inter-model
agreement --- reading agreement as concurrent consistency rather than accuracy,
since such indices likely appear in the models' training data --- and discuss
conceptually how the biases of an LLM annotator and a human expert differ. Code,
prompts, runs, and an interactive gallery of the animated trajectories are
released.
\end{abstract}

\section{Introduction}
\label{sec:intro}

In \emph{The Narrow Corridor}~\citep{acemoglu2019narrow}, Acemoglu and Robinson
propose that liberty is neither the natural product of a strong state nor of a
strong society, but of a contested balance between the two. They picture a
country's history as a path through a plane whose horizontal axis is the power
of civil society and whose vertical axis is the power of the state; a ``narrow
corridor'' running between the axes marks the region where the two powers grow
together and liberty is sustained. The book is rich in narrative but contains
\emph{no plotted trajectories}. The obstacle is not conceptual but
operational: turning centuries of a nation's history into a sequence of
$(\text{society}, \text{state})$ coordinates requires an enormous amount of
expert judgment, and any such coding is vulnerable to the biases, selective
readings, and coarse simplifications that historians rightly distrust.

This is a measurement problem, and measurement problems in the social sciences
increasingly admit machine assistance. Large language models have been shown to
annotate and classify political text at or above the level of crowd
workers~\citep{gilardi2023chatgpt}, and a growing literature examines their use
as instruments in computational social science~\citep{ziems2024can}. We ask a
narrower, concrete question: \emph{can an LLM place a country in the Narrow
Corridor space, period by period, in a way that is reproducible, checkable
against an established expert index, and useful to a researcher?}

We take an applied, digital-humanities stance. Our primary artifact is a
\textbf{six-country trajectory atlas} --- Iran, France, the United Kingdom, the
United States, China, and Chile --- rendered as time-shaded paths through the
corridor, together with the cross-country patterns those paths reveal. The
method and its validation are the backing that makes the atlas trustworthy
rather than decorative. Concretely, we contribute:

\begin{enumerate}
  \item \textbf{A scoring method.} A period-by-period pipeline combining an
        anchored $0$--$10$ rubric, chain-of-thought elicitation of events and
        trends~\citep{wei2022chain}, and in-context anchoring on prior
        periods~\citep{brown2020language,dong2024survey}.
  \item \textbf{A provider-agnostic, open implementation.} Any model reachable
        through a unified interface can be swapped by changing one string; we
        release code, prompts, caches, and every run.
  \item \textbf{A trajectory atlas and cross-country reading} of six countries.
  \item \textbf{An interactive, publicly hosted web gallery.} A static site
        (GitHub Pages) presenting every trajectory as a time-shaded plot with
        \emph{click-to-play animations} that reveal the path period by period
        and annotate the historical event driving each move; readers explore and
        compare models directly in the browser, turning a static figure into a
        navigable atlas.
  \item \textbf{An assessment of the scores} for consistency with an expert
        index (V-Dem) and for inter-model reliability, with explicit treatment of
        the training-data-overlap threat that makes such agreement concurrent
        consistency rather than proof of accuracy.
  \item \textbf{A conceptual discussion of measurement bias}: what an LLM
        annotator plausibly gets right and wrong relative to a human researcher
        or an expert panel doing the identical task (we run no human study here;
        that comparison is left to future work).
\end{enumerate}

\section{Related Work}
\label{sec:related}

\paragraph{The Narrow Corridor and measuring state/society power.}
The (state, society) framing extends the institutional program of Acemoglu and
Robinson~\citep{acemoglu2012why,acemoglu2019narrow}. Quantifying the two axes
connects to decades of effort to measure regime characteristics and state
capacity: the Varieties of Democracy (V-Dem) project aggregates expert codings
into latent indices including a core civil-society index and measures of state
administrative capacity~\citep{coppedge2024vdem,pemstein2018vdem}; Polity5 codes
regime authority on an autocracy--democracy scale~\citep{marshall2020polity};
and Freedom House rates political rights and civil
liberties~\citep{freedomhouse2023}. These projects are themselves large expert
coding operations, which makes them both our natural validation targets
(Section~\ref{sec:validation}) and our natural comparison point when we ask what
an LLM changes about the coding process (Section~\ref{sec:discussion}).

\paragraph{Text as data and LLMs as social-science instruments.}
The use of automated methods to extract measures from text is
well established~\citep{grimmer2013text}. Recent work reports that instruction-tuned
LLMs match or exceed crowd annotators on political text-labeling
tasks~\citep{gilardi2023chatgpt} and surveys their broader promise and pitfalls
for computational social science~\citep{ziems2024can}. A related strand uses LLMs
to \emph{simulate} human populations~\citep{argyle2023out} and debates whether
they can stand in for human participants at
all~\citep{dillion2023can}. Our task differs from labeling or simulation: we ask
the model to produce a continuous, temporally coherent \emph{measurement} of a
theoretical construct, then check that measurement against independent data.

\paragraph{Prompting.}
Our pipeline relies on two now-standard techniques. Chain-of-thought
prompting~\citep{wei2022chain} improves reasoning by eliciting intermediate
steps; we use it to force the model to enumerate historical events and trends
\emph{before} scoring. In-context learning~\citep{brown2020language,dong2024survey}
lets the model condition on examples provided at inference time; we feed all
prior periods' scores forward so the trajectory is temporally coherent rather
than a sequence of independent guesses. This anchoring is a deliberate design
choice with a cost --- an early mis-score can propagate --- which we return to in
Section~\ref{sec:limitations}.

\paragraph{Bias and reliability.}
LLMs inherit biases from training data~\citep{bender2021dangers} and exhibit
measurable political leanings: they do not reflect a representative cross-section
of opinion~\citep{santurkar2023whose}, can carry the political slant of their
pretraining corpora into downstream behavior~\citep{feng2023pretraining}, and
have been probed directly for partisan
bias~\citep{motoki2024more}. These findings motivate both our reliability
analysis --- borrowing the logic of inter-coder reliability from content
analysis~\citep{krippendorff2018content} --- and our discussion of how LLM bias
compares to the well-known biases of human expert coders.

\section{Method}
\label{sec:method}

\paragraph{The space.} We place a country at each time period $t$ at a point
$\big(\mathrm{soc}_t, \mathrm{sta}_t\big) \in [0,10]^2$, where $\mathrm{soc}_t$
is the power of civil society and $\mathrm{sta}_t$ is the power of the state. A
period is a fixed span of years (e.g., a decade). We read a country as
\emph{in the corridor} in period $t$ when its state and society power are close
(near the diagonal) and \emph{above}/\emph{below} it when one clearly dominates.
The dashed lines in our plots mark this band \emph{illustratively}, near the
diagonal, rather than as a fitted threshold, and are drawn at constant width ---
a deliberate simplification of the book's \emph{widening} corridor; positions are
therefore read qualitatively, not from a numeric cutoff.

\paragraph{Anchored rubric.} A free-floating request for ``two numbers'' is not
a measurement. Every scoring prompt embeds a fixed $0$--$10$ rubric that defines
each band of state power (from $0$--$2$, collapsed/nonexistent, to $9$--$10$,
overwhelming penetration of society) and society power (from $0$--$2$, atomized,
to $9$--$10$, society pervasively checks the state), plus explicit neutrality
guidance: judge each period on its own terms, weigh evidence for and against a
shift, avoid presentism, and treat the two axes as independent. The rubric is
what makes scores comparable across periods, countries, and models.

\paragraph{Chain-of-thought: events then scores.} For each period the model is
first asked to list the major events and trends affecting state or society power,
separating slow-moving trends from discrete events and noting the direction of
each effect. This narrative is then inserted into a second prompt that requests
the numeric score, so the number is conditioned on explicit reasoning rather
than produced in one leap~\citep{wei2022chain}.

\paragraph{In-context learning: temporal anchoring.} Starting from the earliest
period (scored on absolute terms), we accumulate every prior period's
$(\mathrm{soc},\mathrm{sta})$ pair and feed the running trajectory into each
subsequent scoring prompt~\citep{brown2020language}. The model reports both the
change from the previous period and the resulting absolute values, which
discourages the sequence from drifting and keeps a quiet period near zero
change.

\paragraph{Structured, validated output.} Each scoring call returns a JSON
object --- the two power values, their per-period changes, a short key event, and
a one-sentence justification --- validated against a fixed schema. A malformed
response triggers one reminder retry and then a numeric fallback extractor before
an error is raised.

\paragraph{Provider-agnostic access and caching.} All model calls go through a
single unified interface, so a model is selected by one identifier string and
providers can be swapped without code changes. Because the pipeline makes one or
two calls per period and is therefore expensive, every response is cached on
disk keyed by (model, prompt); re-runs and visualization iteration are free, and
the released cache makes our exact runs reproducible.

\paragraph{Visualization.} A run renders as (i) a static plot of society vs.\
state power with the path time-shaded from early (dark) to late (bright) and a
year colorbar, the largest moves annotated with their year and key event; and
(ii) an animation that reveals the path period by period, drawing an arrow for
the move each event caused and captioning it. Figure~\ref{fig:atlas} is built
from the static renderings; the animations are best viewed in the interactive
gallery at \url{https://esonghori.github.io/visualize_narrow_corridor_with_ai/}.

\section{Experimental Setup}
\label{sec:setup}

\paragraph{Countries.} We study six countries chosen to span regions of the
corridor and of the world: \textbf{Iran (Persia)}, \textbf{France}, the
\textbf{United Kingdom}, the \textbf{United States}, \textbf{China}, and
\textbf{Chile}. Each is scored from a country-appropriate start year through 2020
(a common cutoff for consistency across cases) in decade (10-year) periods,
giving 14--24 periods per country
(Table~\ref{tab:validation} lists windows and counts). Six countries is a proof
of concept, not a basis for cross-national generalization.

\paragraph{Models.} We compare four models from four families (Google,
Anthropic, OpenAI, Alibaba): three closed-weight models
(\texttt{gemini-3.5-flash}, \texttt{claude-opus-4-8}, \texttt{gpt-4o}) and one
open-weight model (\texttt{qwen-2.5-72b-instruct}). All are driven through the
same interface with identical prompts, isolating model effects from prompt
effects.

\paragraph{Consistency check against an expert index.} We compare the LLM scores
to V-Dem~\citep{coppedge2024vdem}, aligning each period's midpoint year to the
V-Dem year. We map LLM \emph{society power} to the V-Dem core civil-society index
(\texttt{v2xcs\_ccsi}) and LLM \emph{state power} to a state-authority measure
(V-Dem state authority over territory, \texttt{v2terr}) --- a narrow,
territorial-control proxy for the rubric's broader penetration-of-society notion
of state power. The mapping is fixed before inspecting correlations, and
alternative capacity indicators are easy to swap in the released code. We report
rank (Spearman) correlations per country on both \emph{levels} and \emph{first
differences} (period-to-period change): levels are dominated by a shared secular
trend and inflate correlation, so first differences are the more honest
agreement signal. Crucially, V-Dem is
itself almost certainly in the models' training data, so we read this as
\emph{concurrent consistency with expert consensus}, not as external proof of
accuracy (Section~\ref{sec:limitations}). Polity5 and Freedom House are natural
additional anchors but are out of scope here.

\section{Results}
\label{sec:results}

\subsection{Trajectory atlas}
\label{sec:atlas}
Figure~\ref{fig:atlas} shows the six trajectories under \texttt{claude-opus-4-8},
chosen as a single reference model for legibility; the per-model panels for all
four models are in the online gallery, and their disagreement is quantified in
Section~\ref{sec:reliability}. \ph{One paragraph reading the atlas: which
countries sit inside vs.\ above vs.\ below the corridor, and the shape of the
paths (e.g., loops indicating oscillation, long diagonal runs indicating
co-growth).}

\begin{figure*}[t]
  \centering
  % Panels match run_experiments.py slugs: <country-slug>__<model-slug>
  \atlas{iran-persia__anthropic-claude-opus-4-8}\hfill
  \atlas{france__anthropic-claude-opus-4-8}\hfill
  \atlas{united-kingdom__anthropic-claude-opus-4-8}\\[4pt]
  \atlas{united-states__anthropic-claude-opus-4-8}\hfill
  \atlas{china__anthropic-claude-opus-4-8}\hfill
  \atlas{chile__anthropic-claude-opus-4-8}
  \caption{Trajectory atlas: society power (x) vs.\ state power (y) for six
  countries under \texttt{claude-opus-4-8}. Paths are time-shaded (dark $=$
  early, bright $=$ late); dashed lines mark the corridor illustratively.
  Animated versions are available in the online gallery. \ph{Regenerate with
  \texttt{paper/experiments/run\_experiments.py}.}}
  \label{fig:atlas}
\end{figure*}

\subsection{Cross-country patterns}
\ph{Report patterns visible across the atlas: e.g., whether corridor-resident
countries (candidate: UK, France late) show tighter co-growth than
outside-corridor countries (candidate: China high-state/low-society; pre-modern
cases low-state), and how sharp transitions (revolutions, coups) appear as
large single-period jumps.}

\subsection{Consistency with V-Dem}
\label{sec:validation}
Table~\ref{tab:validation} reports correlations between the LLM scores and V-Dem
over overlapping years, on both levels and first differences. \ph{Summarize:
society vs.\ \texttt{v2xcs\_ccsi} and state vs.\ \texttt{v2terr}; contrast the
(high) level correlations against the (lower, more informative) first-difference
correlations; note where consistency holds vs.\ breaks down.} As noted above,
this measures agreement with expert consensus the models were likely trained on,
not independent accuracy.
\begin{table}[t]
\centering
\small
\caption{Consistency with V-Dem (\texttt{claude-opus-4-8}): Spearman $\rho$ at period-midpoint years, reported as level\,/\,$\Delta$ (levels vs.\ first differences). Society vs.\ \texttt{v2xcs\_ccsi}; state vs.\ \texttt{v2svstterr}. First differences are the more honest signal.}
\label{tab:validation}
\begin{tabular}{lcccc}
\toprule
Country & Window & $N$ & Society $\rho$ & State $\rho$ \\
\midrule
Iran (Persia) & 1880--2020 & 14 & 0.45\,/\,0.60 & 0.42\,/\,0.13 \\
France & 1789--2020 & 24 & 0.87\,/\,0.40 & 0.20\,/\,0.04 \\
United Kingdom & 1789--2020 & 24 & 0.70\,/\,0.21 & 0.13\,/\,0.07 \\
United States & 1789--2020 & 24 & -0.10\,/\,0.26 & 0.84\,/\,0.15 \\
China & 1880--2020 & 14 & 0.89\,/\,0.66 & 0.74\,/\,0.24 \\
Chile & 1818--2020 & 21 & 0.74\,/\,0.17 & 0.18\,/\,0.42 \\
\bottomrule
\end{tabular}
\end{table}


\subsection{Model comparison and reliability}
\label{sec:reliability}
Treating each model as a ``rater'' in the sense of inter-coder
reliability~\citep{krippendorff2018content}, Table~\ref{tab:intermodel} reports
agreement between models on the same country. We report Krippendorff's $\alpha$
(interval, treating the 0--10 scores as interval-scaled and the four models as
raters) on the period-to-period \emph{changes} rather than on levels (for the
same reason first differences are used in Section~\ref{sec:validation}):
agreement on changes tests whether models see the same \emph{moves}, not merely
the same secular trend. \ph{Summarize: which axis is more reliable, which model pairs agree
most, and whether the open-weight model diverges from the closed frontier
models.}
\begin{table}[t]
\centering
\small
\caption{Inter-model reliability: Krippendorff's $\alpha$ (interval) on period-to-period \emph{changes} across models, per country. Higher $=$ models see the same moves.}
\label{tab:intermodel}
\begin{tabular}{lccc}
\toprule
Country & \#models & Society $\alpha$ & State $\alpha$ \\
\midrule
Iran (Persia) & 4 & 0.77 & 0.78 \\
France & 4 & 0.72 & 0.56 \\
United Kingdom & 4 & 0.56 & 0.63 \\
United States & 4 & 0.47 & 0.62 \\
China & 4 & 0.56 & 0.56 \\
Chile & 4 & 0.83 & 0.57 \\
Colombia & 4 & 0.54 & 0.53 \\
DR Congo & 4 & 0.51 & 0.77 \\
Lebanon & 4 & 0.34 & 0.71 \\
Zambia & 4 & 0.61 & 0.82 \\
Somalia & 4 & 0.56 & 0.83 \\
India & 4 & 0.47 & 0.16 \\
\midrule
\textbf{Overall} & -- & \textbf{0.66} & \textbf{0.68} \\
\bottomrule
\end{tabular}
\end{table}


\subsection{Qualitative case studies}
\ph{Two or three short reads checking that trajectories match documented history:
Iran around 1979 (sharp state consolidation / society suppression after the
revolution); Chile 1973 and 1990 (coup and redemocratization); and a third
(e.g., UK gradual co-growth, or China post-1949). Note where the model matches
consensus and where it is anachronistic or flattens contested episodes.}

\section{Discussion: LLM vs.\ Human Researcher or Expert Panel}
\label{sec:discussion}

Both an LLM and a human expert (or panel) performing this coding are biased; the
useful question is \emph{how the bias profiles differ}, and which task
properties favor which coder.

\paragraph{What a human panel brings.} Domain expertise, sensitivity to
contested historiography, the ability to flag ``this case is genuinely
disputed,'' and accountable, attributable judgment. Established indices such as
V-Dem manage individual coder bias with multiple experts per case and a latent
measurement model that estimates and corrects for coder
disagreement~\citep{pemstein2018vdem}; content analysis formalizes this as
inter-coder reliability~\citep{krippendorff2018content}. The costs are equally
real: expert panels are slow, expensive, hard to scale to arbitrary
country-periods, and carry their own systematic biases --- national, ideological,
disciplinary --- that are difficult to audit because they live in individual
judgment.

\paragraph{What an LLM changes.} An LLM annotator is fast, cheap at scale,
and \emph{transparent in its instructions}: the rubric and prompts are explicit
artifacts that can be inspected, criticized, versioned, and re-run identically.
Its bias, however, is \emph{opaque in origin}: it reflects the distribution of
its training data~\citep{bender2021dangers,santurkar2023whose,feng2023pretraining},
which over-represents English-language, recent, and Western sources, and it tends
toward received consensus, likely flattening contested or under-documented
episodes rather than flagging them. It has no accountable author, and its scores
can shift with model version or provider (Section~\ref{sec:reliability}).

\paragraph{Reliability as the bridge.} The inter-coder logic that human panels
use for humans applies to models: agreement across independent models is a
measurable proxy for reliability. Where independent models agree on the moves, a
score is more likely to reflect a shared signal than one model's idiosyncrasy;
where they disagree, the disagreement usefully localizes the contested cases a
human panel would also flag. Reliability is necessary but not sufficient for
validity, though --- models trained on overlapping corpora can agree on a shared
bias --- which is why we treat both inter-model agreement and V-Dem consistency
as evidence about \emph{reliability and consensus}, not accuracy.

\paragraph{When to use which.} We argue the LLM pipeline is best positioned as a
\emph{scalable first pass and a reproducible baseline} --- generating candidate
trajectories over many countries cheaply, surfacing where models disagree, and
freeing expert attention for the contested cases --- rather than as a replacement
for expert coding where accountability and contested-case judgment are
paramount.

\section{Limitations}
\label{sec:limitations}
Several limitations bound what this proof of concept establishes.
\emph{Circularity of the consistency check.} V-Dem and similar indices are
almost certainly in the models' training data, so agreement with them cannot
distinguish measuring the construct from recalling memorized values; we therefore
frame the check as concurrent consistency, not accuracy, and disentangling the
two (e.g., counterfactual-history probes, or ablating the events step) is future
work. \emph{Construct compression.} The two axes reduce vast, contested
constructs to single numbers, and our single state-capacity indicator only
partially aligns with the theory's notion of state power. \emph{Coverage and
language bias.} Training corpora over-represent English-language, recent, and
Western sources --- weakest for exactly three of our cases (Iran, China, Chile),
and we prompt in English throughout. \emph{Path dependence.} In-context anchoring
can propagate an early mis-score through the whole trajectory; we do not ablate
it here. \emph{Single-model atlas.} Figure~\ref{fig:atlas} uses one model for
legibility; readers should treat it as one reading, with cross-model spread in
Section~\ref{sec:reliability}. \emph{Scale and statistics.} Six countries and
short per-country series preclude cross-national generalization or strong
significance claims. We deliberately keep the study small; the natural next steps
--- a human expert-coder baseline, prompt/period-granularity ablations, a
non-English prompting probe, more indices, and an ensemble atlas --- are left to
future work.

\section{Conclusion}
We presented, as a proof of concept, a method and a provider-agnostic tool that
operationalize the Narrow Corridor's (state, society) space into reproducible
trajectories, together with a protocol for checking their consistency with
expert consensus and their reliability across models, and released it as a
one-command exercise. The resulting six-country atlas is a first
quantified rendering of a framework that has been purely qualitative, and the
accompanying consistency check and bias discussion frame honestly what such
LLM-generated measurements are --- and are not --- good for.

\section*{Reproducibility}
All code, prompts, the on-disk response cache, and every run (JSON, static
plots, animations) are released at
\url{https://github.com/esonghori/visualize_narrow_corridor_with_ai}. The exact
experiments in this paper are reproduced by
\texttt{paper/experiments/run\_experiments.py} and \texttt{analysis.py}; see
\texttt{paper/README.md}. The interactive gallery of animated trajectories is
hosted at \url{https://esonghori.github.io/visualize_narrow_corridor_with_ai/}.

\bibliographystyle{plainnat}
\bibliography{references}

\end{document}
